% from ega2-1-fr.tex
\oldpage[II]{44}
si $g'$ est un second élément homogène de $S'_+$ et $g=\varphi(g')$, le
diagramme
\[
  \xymatrix{
    (M_{[\varphi]})_{(f')}\ar[r]^{\sim}\ar[d]
    & (M_{(f)})_{[\varphi_{(f)}]}\ar[d]\\
    (M_{[\varphi]})_{(f'g')}\ar[r]^{\sim}
    & (M_{(fg)})_{[\varphi_{(fg)}]}
  }
\]
est commutatif, d'où on conclut aussitôt que l'isomorphisme
\[
  \widetilde{M_{[\varphi]}}|D_+(f'g')
  \xrightarrow{\sim}(\Phi_{fg})_*(\widetilde{M}|D_+(fg))
\]
est la restriction à $D_+(f'g')$ de l'isomorphisme
$\widetilde{M_{[\varphi]}}|D_+(f')
\xrightarrow{\sim}(\Phi_f)_*(\widetilde{M}|D_+(f))$. Comme $\Phi_f$ est la
restriction à $D_+(f)$ du morphisme $\Phi$, on voit donc, compte tenu de
(\hyperref[II.2.8.1.1-fr]{2.8.1.1}), et en posant
$X'=\operatorname{Proj}(S')$~:
\end{env}

\begin{proposition}[2.8.7]
\label{II.2.8.7-fr}
Il existe un isomorphisme canonique fonctoriel du $\mathcal{O}_{X'}$-Module
$\widetilde{M_{[\varphi]}}$ sur le
$\mathcal{O}_{X'}$-Module $\Phi_*(\widetilde{M}|G(\varphi))$.
\end{proposition}

On en déduit aussitôt une application canonique fonctorielle de l'ensemble
des $\varphi$-morphismes $M'\to M$ d'un $S'$-module gradué dans le $S$-module
gradué $M$, dans l'ensemble des $\Phi$-morphismes
$\widetilde{M'}\to\widetilde{M}|G(\varphi)$. Avec les notations de
(\hyperref[II.2.8.4-fr]{2.8.4}), si $M''$ est un $S''$-module gradué, au
composé d'un $\varphi$-morphisme $M'\to M$ et d'un $\varphi'$-morphisme
$M''\to M'$ correspond canoniquement le composé de
$\widetilde{M'}|G(\varphi')\to\widetilde{M}|G(\varphi'')$ et de
$\widetilde{M''}\to\widetilde{M'}|G(\varphi')$.

\begin{proposition}[2.8.8]
\label{II.2.8.8-fr}
Sous les hypothèses de (\hyperref[II.2.8.1-fr]{2.8.1}), soit $M'$ un
$S'$-module gradué. Il existe un homomorphisme canonique fonctoriel $\nu$ du
$(\mathcal{O}_X|G(\varphi))$-Module $\Phi^*(\widetilde{M'})$ dans le
$(\mathcal{O}_X|G(\varphi))$-Module
$\widetilde{M'\otimes_{S'}S}|G(\varphi)$. Si l'idéal $S'_+$
est engendré par $S'_1$, $\nu$ est un isomorphisme.
\end{proposition}

\begin{proof}
En effet, pour $f'\in S'_d$ ($d>0$), on définit un homomorphisme canonique
fonctoriel de $S_{(f)}$-modules (où $f=\varphi(f')$)
\[
\label{II.2.8.8.1-fr}
  \nu_f:M'_{(f')}\otimes_{S'_{(f')}}S_{(f)}
  \longrightarrow(M'\otimes_{S'}S)_{(f)}
\tag{2.8.8.1}
\]
en composant l'homomorphisme
$M'_{(f')}\otimes_{S'_{(f')}}S_{(f)}\to
M'_{f'}\otimes_{S'_{f'}}S_f$ et l'isomorphisme canonique
$M'_{f'}\otimes_{S'_{f'}}S_f\xrightarrow{\sim}(M'\otimes_{S'}S)_f$
(\hyperref[0.1.5.4-fr]{0, 1.5.4}) et notant que ce dernier conserve les
degrés. On vérifie aussitôt la compatibilité des $\nu_f$ avec les opérateurs
de restriction de $D_+(f)$ à $D_+(fg)$ (pour un second $g'\in S'_+$ et
$g=\varphi(g')$), d'où la définition de l'homomorphisme
\[
  \nu:\Phi^*(\widetilde{M'})
  \longrightarrow\widetilde{M'\otimes_{S'}S}|G(\varphi)
\]
compte tenu de (\hyperref[I.1.6.5-fr]{I, 1.6.5}). Pour prouver la seconde
assertion, il suffit de montrer que $\nu_f$ est un isomorphisme pour tout
$f'\in S'_1$, puisque $G(\varphi)$ est alors réunion des
$D_+(\varphi(f'))$. On définit d'abord une application $\mathbf{Z}$-bilinéaire
$M'_m\times S_n\to M'_{(f')}\otimes_{S'_{(f')}}S_{(f)}$ en faisant
correspondre à $(x',s)$ l'élément
$(x'/{f'}^m)\otimes(s/f^n)$ (avec la convention que $x'/{f'}^m$ est
${f'}^{-m}x'/1$

